\section{Definitions, pinning, and the modular triple}

\begin{statement}{A.1}{Proved}{The explicit matrix dictionary}
In the conventions of this paper, the following are functorial
definitions over every commutative ring. Their identities are
proved in A.2--A.9 and Part B. The principal source passages
are C3-259, C3-264--282, C3-321--359, C3-370, C3-393,
C3-399, and C3-479--480.

$$
\begin{array}{c|l}
\text{symbol}&\text{definition}\\ \hline
\rho(A)&rAr^{-1}\\
\sigma(A)&sAs^{-1}\\
\xi&\alpha\rho(\alpha)^{-1}
       =\alpha r\alpha^{-1}r^{-1}\\
\eta&\beta\sigma(\beta)^{-1}
       =\beta s\beta^{-1}s^{-1}\\
\epsilon&-I\\
\epsilon_1&rs=u_1(1)\\
\epsilon_0&-s^{-1}r=sr=u_0(1)\\
T_r&\operatorname{Cent}_{\GL_2}(r)\\
T_s&\operatorname{Cent}_{\GL_2}(s)\\
ST_r&\ker(N_r:T_r\longrightarrow\Gm)\\
ST_s&\ker(N_s:T_s\longrightarrow\Gm)\\
N'_s&\{y:ys=e\,sy,\ e^2=1\}\\
B_0&\left\{\left(\begin{smallmatrix}a&b\\0&d\end{smallmatrix}\right):
                  a,d\text{ units}\right\}\\
B'_0&\left\{\left(\begin{smallmatrix}m&n\\0&1\end{smallmatrix}\right):
                  m\text{ a unit}\right\}\\
U_0&\{u_0(z):z\in R\}\\
U_1&\{u_1(z):z\in R\}\\
M&\{(x,y,f):\det f=\det y=t\det x,\ t^2=1\}\\
M^{00}&\ker(e,t)\\
M^+&\ker(\det f)\\
K&ST_r\times ST_s\\
R_M&\ker(M\longrightarrow\PGL_2)\\
S'T_r&N_r^{-1}(\mu_2)\\
SM&M\times_{\det f,\Gm,\det}B'_0 .
\end{array}
$$

The indices on the two root elements are fixed by their displayed
matrices; a source using the opposite sigma convention must also
translate these indices and signs. In particular $\eta$ is not
$\beta s\beta^{-1}$: the latter is $\eta s$.
\end{statement}
\begin{proof}
The conjugation formulas are definitions. Multiplication of the
fixed matrices gives the three displayed root and central elements.
Commuting a general matrix with $r$ gives
$\left(\begin{smallmatrix}a&b\\-b&a+b\end{smallmatrix}\right)$;
commuting it with $s$ gives
$\left(\begin{smallmatrix}c&d\\-d&c\end{smallmatrix}\right)$.
The relation $ys=e\,sy$ gives the stated signed normalizer in
A.5. The remaining entries are kernels or fibre products of
the explicitly indicated morphisms. Their representability,
dimensions, and exactness properties are proved in Part B.
No equality between the last fibre product and $M^+$ is implied.
\end{proof}

\begin{statement}{A.2}{Proved}{Regular root elements and their fixed lines}
Let $g\in\SL_2(R)$ have trace two. Then $N=g-I$ satisfies $N^2=0$.
If the entries of $N$ generate the unit ideal, its image and kernel
are the same locally direct-summand line. Locally on the base,
$g$ is conjugate to
$\left(\begin{smallmatrix}1&1\\0&1\end{smallmatrix}\right)$.
Its unique fixed line is $\ker N$, and the associated root subgroup
is canonically the vector group
$\operatorname{Hom}(R^2/\ker N,\ker N)$.
These statements concern the indicated primitive root elements;
they do not identify all abstract unipotent elements over
nonreduced rings.
\end{statement}
\begin{proof}
Cayley--Hamilton gives $g^2-2g+I=0$, hence $N^2=0$.
Also $\det N=\det g-\tr g+1=0$.
On an open set where an entry of $N$ is a unit, elementary row
and column operations exhibit a rank-one map with projective
image and kernel. The inclusion $\im N\subseteq\ker N$ is an
equality because both are direct-summand lines and the induced
map from the quotient to the image is an isomorphism.
Choose a basis of the line and a quotient basis so that this
isomorphism has coefficient one. This gives the displayed matrix.
A fixed line on which the induced root action is trivial must
lie in $\ker N$, hence equals it. Maps from the quotient line
to the fixed line square to zero; $h\mapsto I+h$ therefore
identifies their additive vector group with the root subgroup.
This is the rank-two instance of the root-group conventions
in \citet[Expos\'e XXIII, Definitions1.1--1.2]{AB_SGA3};
the argument above proves the particular integral assertion used.
\end{proof}

\begin{statement}{A.3}{Proved}{Canonical algebraic pinning}
Let $A$ be an Azumaya algebra of rank four over a ring $R$.
Suppose $r_0,s_0\in A^\times$ satisfy

$$
\det r_0=\det s_0,\qquad
 \det(r_0+s_0)=0,\qquad
 u=\tr(r_0+s_0)\in R^\times .
$$

Reduced trace and determinant are meant. Then there is a unique
algebra isomorphism with $M_2(R)$ under which

$$
r_0=\begin{pmatrix}0&u\\v&w\end{pmatrix},\qquad
 s_0=\begin{pmatrix}0&-u\\-v&u-w\end{pmatrix},
 \quad v=-\det(r_0)/u\in R^\times,\quad w=\tr(r_0).
$$

Thus these hypotheses themselves split the Azumaya algebra.
The uniqueness is uniqueness of this algebraic pinning,
not uniqueness among all abstract group automorphisms of $\GL_2$.
\end{statement}
\begin{proof}
Put $p=(r_0+s_0)/u$. Cayley--Hamilton gives $p^2=p$,
$\tr p=1$, and $\det p=0$. After a faithfully flat splitting
of $A$, this is a rank-one projector. Write $q=1-p$.
In a local basis with $p=E_{22}$, write
$r_0=\left(\begin{smallmatrix}a&b\\c&d\end{smallmatrix}\right)$.
Then
$s_0=\left(\begin{smallmatrix}-a&-b\\-c&u-d\end{smallmatrix}\right)$,
and
$\det s_0-\det r_0=-ua$.
Consequently $a=0$ and $bc=-\det r_0$ is a unit.

The following globally defined elements therefore form matrix units:

$$
E_{22}=p,\quad E_{11}=q,\quad
 E_{12}=u^{-1}q r_0p,\quad
 E_{21}=-u(\det r_0)^{-1}p r_0q.
$$

Their products can be checked after the faithfully flat splitting,
where $E_{12}E_{21}=q$ and $E_{21}E_{12}=p$.
They give an algebra isomorphism $M_2(R)\to A$.
The coefficients of $r_0,s_0$ in these units are exactly those
displayed. Every compatible isomorphism must preserve $p,q$
and the two off-diagonal units, so is identical.
\end{proof}

\begin{statement}{A.4}{Refuted}{Uniqueness of a projective unipotent lift}
A primitive determinant-one, trace-two lift of a given projective
root element need not be unique. The scalar ambiguity is precisely

$$
g\longmapsto(1+\tau)g,\qquad 2\tau=0,\quad \tau^2=0.
$$

On a fixed base it vanishes exactly when there is no such nonzero
global scalar. It vanishes after every base change exactly when
two is invertible. Rings of cardinality four already give
counterexamples.
\end{statement}
\begin{proof}
Two lifts to $\GL_2$ of the same projective point differ by a
unique scalar unit $\lambda$. Requiring determinant one and
trace two gives $\lambda^2=1$ and $2\lambda=2$.
Writing $\lambda=1+\tau$ gives the two equations stated.
Conversely they suffice; primitivity is preserved.

Over $\Z/4$, the distinct matrices

$$
\begin{pmatrix}1&0\\1&1\end{pmatrix},\qquad
 \begin{pmatrix}3&0\\3&3\end{pmatrix}
$$

have determinant one, trace two, square-zero differences from
the identity, and the same projective image.
The lower off-diagonal entries are units.
No field has a nonzero square-zero scalar, so no smaller finite
ring is a counterexample. If two is not invertible, base change
to $(R/2R)[\tau]/(\tau^2)$ produces a nonzero ambiguity.
If two is invertible, $2\tau=0$ forces $\tau=0$.
\end{proof}

\begin{statement}{A.5}{Proved}{Full centralizers and the signed normalizer}
The full centralizers have dimensions two:

$$
T_r=\Res_{\Z[T]/(T^2-T+1)/\Z}\Gm,\qquad
 T_s=\Res_{\Z[i]/\Z}\Gm.
$$

Their determinants are $n=a^2+ab+b^2$ and $q=c^2+d^2$.
They are not the one-dimensional norm-one groups $ST_r,ST_s$.
Moreover

$$
N'_s=\left\{
 y(c,d,e)=\begin{pmatrix}c&d\\-ed&ec\end{pmatrix}:
 e^2=1,\ q\text{ a unit}\right\},
\quad \det y=eq .
$$

The sign is uniquely determined by $y$, even over nonreduced
rings. The multiplication law is

$$
(c,d,e)(c',d',e')
 =(cc'-e'dd',\,cd'+e'dc',\,ee').
$$

\end{statement}
\begin{proof}
The commutant calculations in A.1 give the two quadratic
algebras and their unit loci. Direct multiplication gives their
norms. Solving $ys=e\,sy$ gives the displayed matrix.
If two signs work, their difference times both $c,d$ is zero;
since $c^2+d^2$ is a unit, the signs agree.
The multiplication formula follows by multiplying the matrices.
Thus the signed normalizer is a semidirect product of $T_s$
with $\mu_2$, with the indicated conjugation action.
The quadratic algebras fail to be etale at three and two,
respectively. Calling their full unit groups norm-one tori
changes both their dimension and their functor of points.
\end{proof}

\begin{statement}{A.6}{Proved}{The universal Borel commutator chart}
For
$U=\left(\begin{smallmatrix}m&n\\0&1\end{smallmatrix}\right)$,
put $D=m^{-1}$ and $B=n/m$. Then

$$
[U,s]=
 \begin{pmatrix}(1+B^2)/D&B\\B&D\end{pmatrix},
$$


$$
[U,r]=
 \begin{pmatrix}(1+B+B^2)/D-B&B\\1+B-D&D\end{pmatrix}.
$$

Consequently the scheme of pairs of matrices with exactly these
shared entries, determinant one, and $D$ a unit is isomorphic
over $\Z$ to $B'_0\simeq\Gm\times\Ga$.
For a fixed chart point, all solutions of

$$
[\alpha,r]=[U,r],\qquad [\beta,s]=e[U,s],\quad e^2=1
$$

are uniquely

$$
\alpha=Ux^{-1},\qquad \beta=Uy^{-1},
 \quad x\in T_r,\quad y\in N'_s,\quad ys=e\,sy.
$$

This describes the stated chart and its fibres. It is not
an identification of this chart with the full group $M$.
\end{statement}
\begin{proof}
The two commutators are obtained by multiplying $U,r,s$
and their inverses. Conversely the common entries recover
$m=D^{-1}$ and $n=B/D$. The determinant equations recover
the upper-left entries, so no reducedness or field assumption
is used.
The first commutator equation is equivalent to
$\alpha^{-1}U$ commuting with $r$.
The second is equivalent to
$(\beta^{-1}U)s=e\,s(\beta^{-1}U)$.
Thus $x=\alpha^{-1}U$ and $y=\beta^{-1}U$ are the unique
parameters. All these identities are checked over the
universal Laurent polynomial ring by
\texttt{code/ab\_reconstructed.js}.
\end{proof}

\begin{statement}{A.7}{Proved}{Normalized modular triples over the requested bases}
Call a triple normalized if $\epsilon=-I$ and

$$
\det r_0=\det s_0=1,\quad
 \tr r_0=1,\quad\tr s_0=0,\quad
 \det(r_0+s_0)=0.
$$

Over $\Z$, over every field, and over $\Z/4,\Z/8,\Z/9$,
all such triples form one $\GL_2$-conjugacy orbit.
Their $\SL_2$-conjugacy classes are naturally indexed by
$R^\times/R^{\times2}$.
The relations $r_0^3=s_0^2=-I$ follow.
In characteristic two these are relations with $\epsilon=I$;
an additional requirement that $\epsilon$ have exact order two
would instead make that locus empty.
\end{statement}
\begin{proof}
A.3 applies with $u=1,v=-1,w=1$, giving the fixed matrices.
For the listed bases every rank-one projective module is free.
The matrix units in A.3 therefore arise from a basis of $R^2$,
so the algebraic pinning is induced by $\GL_2(R)$ conjugation.
The common centralizer is the scalar group because the two
matrices generate all of $M_2(R)$.
Hence the orbit is $\GL_2(R)/R^\times$ and its $\SL_2$-orbits
are the double quotient

$$
\SL_2(R)\backslash\GL_2(R)/R^\times
 \simeq R^\times/R^{\times2}.
$$

The determinant map is surjective and scalar determinants
are squares, proving the last isomorphism.
Cayley--Hamilton gives $r_0^2-r_0+I=0$ and $s_0^2+I=0$,
hence the asserted relations.
The square-class formulation also covers imperfect fields
of characteristic two; their units need not all be squares.
\end{proof}

\begin{statement}{A.8}{Proved}{The finite orbit census}
For the indicated rings, exhaustive enumeration gives

$$
\begin{array}{c|r|r|r}
R&\text{triples}&\SL_2\text{ orbits}&\text{size of each orbit}\\\hline
\F_2&6&1&6\\
\F_3&24&2&12\\
\F_5&120&2&60\\
\F_7&336&2&168\\
\F_{11}&1320&2&660\\
\F_{13}&2184&2&1092\\
\Z/4&48&2&24\\
\Z/8&384&4&96\\
\Z/9&648&2&324
\end{array}
$$

The signed chart of A.6 has respectively $12,576,1536$
distinct pairs over $\F_2,\F_3,\Z/4$.
\end{statement}
\begin{proof}
A.7 gives the orbit count and orbit size
$|\SL_2(R)|/|\mu_2(R)|$. The replacement program independently
enumerates every determinant-one trace-one matrix and every
determinant-one trace-zero matrix, retaining precisely the
pairs with singular sum. It checks their total against
$|\SL_2(R)|$ and verifies the square-class partition.
For the signed chart it enumerates all $U,x,y$ and stores
the resulting ordered matrix pairs. Equality of the number
of distinct images with the number of parameters verifies
injectivity and the displayed cardinalities on those rings.
The universal inverse in A.6 supplies the all-ring proof.
\end{proof}

\begin{statement}{A.9}{Refuted}{Two invalid weakenings}
The algebraic pinning is not unique among all group-scheme
automorphisms of $\GL_2$. Also regularity of two unrelated
matrices and equality of their twisted commutators do not
suffice for the Borel parametrization of A.6.
\end{statement}
\begin{proof}
The automorphism
$g\mapsto(\det g)^{-1}g$ is an involution of $\GL_2$ fixing
$\SL_2$ pointwise. It fixes $r,s$ but is not the identity,
for example on $\operatorname{diag}(2,1)$ over $\Q$.

For the second assertion work over $\F_5$, set $r'=s'=s$,
and take
$\alpha=\beta=\left(\begin{smallmatrix}0&1\\1&2\end{smallmatrix}\right)$.
The commutators are equal, and $r',s'$ are regular.
But $\alpha$ cannot belong to $B'_0T_s$: the lower row of
a matrix in that product is $(-d,c)$ with $c^2+d^2$ a unit,
whereas the lower row $(1,2)$ has norm $1+4=0$.
This obstruction persists after faithfully flat base change.
It refutes the weakened conditions just stated; it does not
remove the pinning and chart hypotheses in A.3 and A.6.
\end{proof}

\begin{statement}{A.10}{Unfinished}{The complete later symbol-by-symbol concordance}
The explicit core dictionary above and the formal constructions
in B.13--B.22 do not constitute a recovered, complete transcription
of every later symbol in §§48--49. The inaccessible original
working table included further entries and scan concordances.
Those entries cannot now be certified verbatim.
In particular, the source-specific restricted-normalizer
projections around C3-683--695 must not be declared surjective
merely from the formal fibre-product argument.
\end{statement}
\begin{proof}[What is established and what is missing]
B.17 proves the exact reconstruction theorem with its explicit
equivariance and kernel hypotheses. B.18 gives a finite
counterexample to an unrestricted later normality assertion.
B.19--B.22 identify the exact image, cocycle, and diagonal-gluing
conditions required in the subsequent constructions.
The finite counterexample and the chart computations were
executed again in the surviving runtime.
What is missing from this fallback is the full recovered
source-specific list of restricted functors and a proof of
the required image equalities for its designated universal
profinite input. Accordingly the unrestricted dictionary and
all-formulas subrequests are not marked complete.
\end{proof}

\begin{statement}{A.11}{Proved}{A typed dictionary for the later formal constructions}
For group data $(C,L,B)$ with $B\subseteq\Aut(C)$ preserving $L$,
the construction of B.13 is a quotient group, or a profinite
quotient when all data are profinite and closed.
For equivariant data $(C,N,A,\phi)$ as in B.17 it is the
diagonal quotient $(C\rtimes N)/D_\phi$.
For maps to a common quotient $\Gamma$, the constructions in
B.19 and B.21 are fibre products and inverse-image subgroups.
For several marked extensions by $C$, the construction in B.22
is a diagonal subextension of their fibre product.
None of these statements alone asserts representability of an
arbitrary fppf sheaf quotient by an affine group scheme.
\end{statement}
\begin{proof}
The exact hypotheses and proofs are given in the cited statements.
Quotients in groups, closed quotients in profinite groups, and
quotient fppf sheaves have different existence theorems.
The constructions provide the corresponding objects in the
categories explicitly named. Representability, where needed,
is proved separately for the finite-dimensional matrix schemes
of Part B and the affine inverse limits of Part E.
This distinction prevents a formal group calculation from being
used as an unproved representability or point-surjectivity claim.
\end{proof}
